\transition{FADE IN:}
\sceneheading{INT. GHAZNI, BIRUNI'S QUARTERS - DAY}
\action{Biruni back among his own instruments. They read differently now --- competence that used to cost him nothing now sits oddly against two years of effort. He unpacks the notebooks from India: the accumulated mess, the unresolved contradictions Pandit told him not to make clean.}
\action{He lays them out on the table in some rough order. Doesn't start writing yet. Just looks at the volume of it --- crossed-out pages, marginal notes, the second kind of entry from the Pramāṇa scene sitting alongside ordinary translations.}
\transition{CUT TO:}
\sceneheading{INT. GHAZNI COURT, ANTECHAMBER - DAY}
\action{A COURT FIGURE greets Biruni with vague, uncomprehending warmth.}
\charcue{COURT FIGURE}
\begin{dialogue}
They say you learned their whole religion.\\
\end{dialogue}
\charcue{BIRUNI}
\begin{dialogue}
I learned a fraction of it. Carefully.\\
\end{dialogue}
\charcue{COURT FIGURE}
\begin{dialogue}
\parenthetical{already moving on}
Good, good.\\
\end{dialogue}
\transition{CUT TO:}
\sceneheading{INT. GHAZNI, SHARED WORKSPACE - DAY}
\action{The COURT ASTROLOGER, older now, still doing exactly what he did in Act One --- reading the instruments toward whatever answer is wanted. He and Biruni share the space without much conversation. The contrast needs no dialogue: the Astrologer unchanged and comfortable, Biruni visibly changed and visibly tired in a way that isn't about to go away.}
\transition{CUT TO:}
\sceneheading{INT. CLERIC'S CHAMBER - DAY}
\action{The CLERIC from the Ghaznavid camp --- or a superior standing in for him --- requests to see the manuscript-in-progress. Direct pressure now.}
\charcue{CLERIC}
\begin{dialogue}
Nobody is asking you to lie.\\
\end{dialogue}
\charcue{BIRUNI}
\begin{dialogue}
No. You're asking me to leave out the parts that are true and inconvenient. That's a different kind of lying. Not a lesser one.\\
\end{dialogue}
\charcue{CLERIC}
\begin{dialogue}
Consider carefully who reads this, once it's finished. And who decides what you're given to finish it with.\\
\end{dialogue}
\action{He doesn't threaten outright. He doesn't need to. The implication --- patronage, access, material support, all of it contingent --- sits in the room without being named.}
\transition{CUT TO:}
\sceneheading{INT. BIRUNI'S QUARTERS - NIGHT}
\action{Writing. Not a triumphant montage --- a laborious one.}
\action{Draft pages fill with cross-outs. The camera holds on the page in extreme close-up at intervals, echoing --- deliberately, if a match-cut is achievable --- the Sanskrit tablet notes from Khwarezm, then the Pramāṇa page from India, now the manuscript itself. Same hand. Same habit. Increasingly abstract objects.}
\action{At one point, Biruni writes a passage, reads it back, and catches himself doing exactly what Pandit caught in the study scene months earlier --- smoothing something foreign into something too comfortably familiar. This time there is no one else in the room to catch it.}
\action{He crosses it out. Rewrites it harder --- a beat longer, a beat more honest --- without anyone watching him do it.}
\action{Later in the same night, working through material on traditions he encountered only by report rather than direct study, he writes a line acknowledging exactly that --- that for some of what follows, he is relying on what he was told rather than what he verified himself. He underlines it, rather than smoothing it into the surrounding text as if it carried the same weight as his firsthand work.}
\transition{CUT TO:}
\sceneheading{INT. GHAZNI COURT, AUDIENCE HALL - DAY}
\action{Biruni before MAHMUD again. Older. Same transactional register as their first exchange, though something in it has shifted with time rather than with any change in Mahmud himself.}
\charcue{MAHMUD}
\begin{dialogue}
Two years, and what do I have.\\
\end{dialogue}
\charcue{BIRUNI}
\begin{dialogue}
A book that won't tell you what you already believe.\\
\end{dialogue}
\charcue{MAHMUD}
\begin{dialogue}
\end{dialogue}
\charcue{(a beat --- almost respect,}
\begin{dialogue}
\end{dialogue}
\charcue{not quite)}
\begin{dialogue}
That's not nothing. It's also not much.\\
\end{dialogue}
\charcue{BIRUNI}
\begin{dialogue}
It's what's true. I told you that's what I'd bring back.\\
\end{dialogue}
\action{Mahmud doesn't stop him. Doesn't praise him either. Already turning his attention elsewhere --- the same gesture from Act One, not warmed by two years, only repeated. That repetition is the point: the work's value was never going to depend on this room recognizing it.}
\transition{CUT TO:}
\sceneheading{INT. BIRUNI'S QUARTERS - NIGHT}
\action{The preface.}
\action{Shot as close to one unbroken take as the scene allows. Biruni alone, writing. Both notebooks --- India's, and the manuscript's own draft history --- present in the room without being shown directly; their weight is in how he writes, not in cutaways to them.}
\action{He writes, in his own words (paraphrased here, not quoted verbatim from any historical text): that he intends to set down what the people of India hold to be true, in terms they would recognize as their own --- neither defending it nor refuting it, only representing it as it stands.}
\action{He finishes the passage. Sets down the stylus. No one is in the room to witness this. No music swell, no triumphant turn to camera. He looks at the page the way he looked at the number in the dirt on the mountain, at the start of the film: private, settled, no external validation attached to it or needed by it.}
\transition{CUT TO:}
\sceneheading{INT./EXT. GHAZNI, VARIOUS - DAY}
\action{Brief coda. A SCRIBE copies the manuscript by hand --- ordinary, unglamorous, ink and patience. No sense of the book's later importance; just the plain mechanics of how consequential work actually enters the world, one copied page at a time.}
\transition{CUT TO:}
\sceneheading{EXT. GHAZNI OR UNSPECIFIED, LATER - DAY}
\action{Biruni, aged now, doing something small and specific --- measuring, calculating, still at it, still unhurried. No summing-up dialogue. The accumulated craft stands for itself.}
\transition{CUT TO:}
\sceneheading{EXT. MOUNTAINSIDE - PRE-DAWN}
\action{The same location as the opening. The same stillness.}
\action{A figure climbs the last stretch of the slope, carrying an instrument case --- NOT Biruni. Younger. We don't need to know exactly who; it could be the grown Shepherd from the opening, or simply someone new to whom the practice has passed. The film doesn't explain the substitution. It trusts the image.}
\action{The figure sets up a quadrant. Sights toward the horizon. Marks something down. The same gesture, exactly, that opened the film --- except the man who taught it to no one in particular is nowhere in the frame.}
\action{The horizon hasn't moved.}
\transition{FADE OUT.}
\action{END OF ACT THREE. END OF SCREENPLAY.}
